% =====================================================================
%  theory_appendix.tex
%  Drop-in theory section for paper_draft.tex.
%  Adds formal definitions, four theorems, and matching empirical
%  validation hooks so the paper stops being "multi-regime benchmark"
%  and becomes "theory of adversarial unlearning with matching audit."
%
%  To include from paper_draft.tex, insert:
%     % =====================================================================
%  theory_appendix.tex
%  Drop-in theory section for paper_draft.tex.
%  Adds formal definitions, four theorems, and matching empirical
%  validation hooks so the paper stops being "multi-regime benchmark"
%  and becomes "theory of adversarial unlearning with matching audit."
%
%  To include from paper_draft.tex, insert:
%     % =====================================================================
%  theory_appendix.tex
%  Drop-in theory section for paper_draft.tex.
%  Adds formal definitions, four theorems, and matching empirical
%  validation hooks so the paper stops being "multi-regime benchmark"
%  and becomes "theory of adversarial unlearning with matching audit."
%
%  To include from paper_draft.tex, insert:
%     % =====================================================================
%  theory_appendix.tex
%  Drop-in theory section for paper_draft.tex.
%  Adds formal definitions, four theorems, and matching empirical
%  validation hooks so the paper stops being "multi-regime benchmark"
%  and becomes "theory of adversarial unlearning with matching audit."
%
%  To include from paper_draft.tex, insert:
%     \input{theory_appendix}
%  immediately after the "Threat Models and Attack Regimes" section.
% =====================================================================

\section{A Theory of Adversarial Unlearning}
\label{sec:theory}

The empirical regimes of this paper measure a single object: the smallest
input-space perturbation that re-elicits the forgotten class on a per-example
basis. This section names that object, ties it quantitatively to certified
unlearning, separates selective leakage from generic adversarial fragility,
and gives a constructive impossibility result for the open
``local-delivery transfer'' problem. The benchmark sections that follow are
re-cast as empirical validation of these statements rather than as standalone
attack curves.

\subsection{Setup and the residual forget-channel capacity}

\paragraph{Notation.} Let $\mathcal{X} = [0,1]^d$ be the input space and
$\mathcal{Y} = [K]$ the label set. Fix the forget class $t \in [K]$. Let
$D = D_R \cup D_F$ with $D_F = \{(x_i, t)\}_{i=1}^{n_f}$. Let $\mathcal{A}$ be
an unlearning algorithm producing $\theta_u = \mathcal{A}(D, t)$. Let
$\theta_o$ denote a from-scratch oracle trained on $D_R$. Let
$f_\theta : \mathcal{X} \to \mathbb{R}^K$ be the logit map and define the
forget-class \emph{target margin}
\begin{equation}
  m_t(x;\theta) \;:=\; f_\theta(x)_t - \max_{j \neq t} f_\theta(x)_j.
\end{equation}
The unlearned model predicts the forgotten class at $x$ iff
$m_t(x;\theta_u) > 0$.

\paragraph{Recovery radius.} For an $L_\infty$ budget,
\begin{equation}
  r_t(x;\theta) \;:=\; \inf\bigl\{\epsilon \geq 0 : \exists\, x' \in
   B_\infty(x,\epsilon) \cap \mathcal{X},\ m_t(x';\theta) > 0\bigr\},
\end{equation}
with $r_t(x;\theta) = 0$ when the clean prediction already matches. The
per-example audit in \texttt{scripts/audits/17\_recovery\_radius\_audit.py}
returns a finite-budget estimator $\hat r_t(x;\theta_u)$ via probe-then-binary
search with PGD / APGD inner loops.

\paragraph{The right population object.} Let $\mu_t$ denote the
forget-class input distribution and $\mu_{\neg t}$ the matched non-forget
control distribution used throughout the paper. Define the
\emph{residual forget-channel capacity} at scale $\epsilon$:
\begin{equation}
  \boxed{\;\mathcal{R}_\epsilon(\theta_u)
  \;:=\;
  \Pr_{x \sim \mu_t}\!\bigl[r_t(x;\theta_u) \leq \epsilon\bigr]
  \;-\;
  \Pr_{x \sim \mu_{\neg t}}\!\bigl[r_t(x;\theta_u) \leq \epsilon\bigr].\;}
  \label{eq:rfcc}
\end{equation}
$\mathcal{R}_\epsilon \in [-1,1]$ is the signed difference between forget-side
and retain-side recoverability at audit scale $\epsilon$. The forget--retain
gap reported throughout the paper is exactly the empirical estimator of
$\mathcal{R}_\epsilon$. Two qualitative regimes the paper invokes informally
become precise statements about $\mathcal{R}_\epsilon$:
\begin{itemize}[leftmargin=1.5em,topsep=0pt]
  \item \textbf{Selective leakage} (paper's ORBIT/CE-U story):
    $\mathcal{R}_\epsilon(\theta_u) > 0$ — forget inputs are strictly easier
    to recover than retain controls.
  \item \textbf{Generic fragility} (paper's SCRUB story):
    $\mathcal{R}_\epsilon(\theta_u) \approx 0$ — the model is adversarially
    fragile, but not selectively so on the forget class.
\end{itemize}

\subsection{Theorem 1: the recovery--certification gap}
\label{sec:thm1}

We bridge the per-example recovery radius to any \emph{certified-unlearning}
budget the method might claim.

\begin{theorem}[Recovery--certification gap]
\label{thm:gap}
Suppose the target-margin function is $L$-Lipschitz in $\|\cdot\|_\infty$
over $\mathcal{X}$, i.e.\ $|m_t(x;\theta_u) - m_t(x';\theta_u)| \leq
L\,\|x - x'\|_\infty$ for all $x,x' \in \mathcal{X}$. Then for every
$x \in \mathcal{X}$,
\begin{equation}
  r_t(x;\theta_u) \;\geq\; \frac{-m_t(x;\theta_u)}{L}.
  \label{eq:thm1-pointwise}
\end{equation}
Equivalently,
\begin{equation}
  \Pr_{x \sim \mu_t}\!\bigl[r_t(x;\theta_u) \leq \epsilon\bigr]
  \;\leq\;
  \Pr_{x \sim \mu_t}\!\bigl[\,m_t(x;\theta_u) \geq -L\epsilon\,\bigr].
  \label{eq:thm1-population}
\end{equation}
\end{theorem}

\begin{proof}
For any $\epsilon \geq 0$ and any $x' \in B_\infty(x,\epsilon)\cap\mathcal{X}$,
the Lipschitz bound gives
$m_t(x';\theta_u) \leq m_t(x;\theta_u) + L\epsilon$. For the recovery event
$m_t(x';\theta_u) > 0$ to hold we therefore require
$\epsilon > -m_t(x;\theta_u)/L$, which proves \eqref{eq:thm1-pointwise}.
Then \eqref{eq:thm1-population} follows by integrating against $\mu_t$.
\end{proof}

\begin{corollary}[Certified unlearning needs separated margins]
\label{cor:cert}
Let $\theta_u$ be $(\epsilon_c,\delta)$-certified relative to the oracle
$\theta_o$ in the conditional TV sense:
\[
\Pr_{x \sim \mu_t}\!\Bigl[
\mathrm{TV}\bigl(\pi_{\theta_u}(\cdot \mid x),\, \pi_{\theta_o}(\cdot \mid x)\bigr)
\leq \epsilon_c\Bigr] \;\geq\; 1 - \delta,
\]
where $\pi_\theta = \mathrm{softmax}\circ f_\theta$. Assume the oracle has an
expected separation $\mathbb{E}_{\mu_t}\!\bigl[-m_t(x;\theta_o)\bigr] \geq
\gamma > 0$ and the logits are bounded $\|f_\theta\|_\infty \leq B$. Then with
probability $\geq 1-\delta$,
\begin{equation}
  \mathbb{E}_{x \sim \mu_t}\!\bigl[\,r_t(x;\theta_u)\,\bigr]
  \;\geq\; \frac{\gamma - 4B\,\epsilon_c}{L}.
\end{equation}
\end{corollary}

\begin{proof}
For any logit map $f$, the margin $m_t = f_t - \max_{j\neq t} f_j$ is a
$1$-Lipschitz functional of $f$ (the maximum of $K-1$ linear functionals).
By the variational form of TV applied to the bounded functional $m_t$ on the
probability simplex $\Delta^{K-1}$,
\[
|m_t(x;\theta_u) - m_t(x;\theta_o)| \;\leq\; 4B \cdot
\mathrm{TV}(\pi_{\theta_u}(\cdot\mid x),\pi_{\theta_o}(\cdot\mid x)).
\]
Take expectations under $\mu_t$, use the high-probability TV bound, then
apply Theorem~\ref{thm:gap}.
\end{proof}

\paragraph{What this says about the benchmark.}
Every method in the paper has a measured $\hat r$ distribution and a
straightforwardly estimable Lipschitz proxy $\hat L$ (e.g.\ the spectral
norm of the Jacobian of the margin at audit points; the
\texttt{16\_jacobian\_access\_audit.py} machinery already computes the
underlying quantity). Corollary~\ref{cor:cert} contrapositively gives a
\emph{provable upper bound on the certifiable unlearning budget}:
\[
  \epsilon_c \;\leq\; \frac{\gamma - L\cdot \mathrm{median}(\hat r)}{4B}.
\]
A method whose median recovery radius is small relative to the oracle margin
\emph{cannot} be honestly marketed at small $\epsilon_c$; the bound is a
falsification test that audits with no access to the oracle distribution can
already perform. This converts the paper's empirical numbers into a stake in
the certified-unlearning literature
\cite{chien2024langevin,hsu2025rurk,yu2025impossibility}.

\subsection{Theorem 2: selective leakage is a Pinsker test}
\label{sec:thm2}

\begin{theorem}[Pinsker bound on selective leakage]
\label{thm:pinsker}
Let $p_F := \Pr_{\mu_t}[r_t \leq \epsilon]$ and
$p_R := \Pr_{\mu_{\neg t}}[r_t \leq \epsilon]$. Then
\begin{equation}
  |\mathcal{R}_\epsilon(\theta_u)| \;=\; |p_F - p_R|
  \;\leq\; \sqrt{\tfrac{1}{2}\,
   \mathrm{KL}\bigl(\mathrm{Ber}(p_F)\,\|\,\mathrm{Ber}(p_R)\bigr)}.
\end{equation}
Consequently, $|\mathcal{R}_\epsilon(\theta_u)| \geq c$ forces
$\mathrm{KL}\bigl(\mathrm{Ber}(p_F)\|\mathrm{Ber}(p_R)\bigr) \geq 2c^2$,
i.e.\ the audit's binary recoverability outcome distinguishes forget from
retain at variational rate $c$.
\end{theorem}

\begin{proof}
Pinsker's inequality applied to the two Bernoulli laws on the binary outcome
$\mathbf{1}\{r_t \leq \epsilon\}$.
\end{proof}

\begin{corollary}[SCRUB is a generic-fragility regime]
\label{cor:scrub}
If $\mathcal{R}_\epsilon(\theta_u) \to 0$ as $n_f \to \infty$ at some
audit scale $\epsilon$, no unlearning claim about $\mu_t$-specific
information removal is empirically supported at that scale: the
data-conditional binary leakage event is indistinguishable from the
retain-control baseline.
\end{corollary}

\paragraph{Why this matters.} The paper's qualitative observation that SCRUB
``exhibits generic adversarial fragility rather than selective forgetting
leakage'' becomes a hypothesis test: the binary recoverability outcomes on
forget and retain controls are samples from $\mathrm{Ber}(p_F)$ and
$\mathrm{Ber}(p_R)$; Theorem~\ref{thm:pinsker} bounds the variational
information they carry about the forget condition. The same machinery
sharpens the ORBIT/CE-U positive results: the gap is not just bigger, it
\emph{forces} a quantitative KL.

\subsection{Theorem 3: local-delivery transfer impossibility}
\label{sec:thm3}

This formalizes the paper's open problem (``patch-aware defenses fail to
transfer to frame or multi-patch delivery'') as a structural impossibility.

\paragraph{Delivery families and covering.} A \emph{delivery mask} is a
subset $M \subseteq [d]$ of input coordinates. A \emph{delivery family} is a
collection $\mathcal{D} \subseteq 2^{[d]}$. A delivery-conditional
perturbation under $\mathcal{D}$ with budget $\epsilon$ is a
$\delta \in \mathbb{R}^d$ with $\mathrm{supp}(\delta) \subseteq M$ for some
$M \in \mathcal{D}$ and $\|\delta\|_\infty \leq \epsilon$. We say a finite
collection of families $\{\mathcal{D}_i\}_{i=1}^k$ is \emph{covering at scale
$\epsilon$} on $\mathcal{X}$ if for every $x_1, x_2 \in \mathcal{X}$ there is
a sequence $x_1 = y_0, y_1, \dots, y_m = x_2$ with each $y_{j+1} - y_j$ a
valid delivery-conditional perturbation under some $\mathcal{D}_{i(j)}$ and
budget $\epsilon$.

\begin{lemma}[Coverage by patches plus frames]
\label{lem:cover}
For $d_{\rm img} = 224 \times 224$ images normalized to $[0,1]^d$, the union
of (i) all $32\times 32$ square delivery masks at every position and (ii)
all $16$-pixel border frames is covering at scale $\epsilon = 1$.
\end{lemma}

\begin{proof}
Any pixel of the image is covered by at least one $32 \times 32$ square or
lies on the border frame. With budget $\epsilon = 1$, a single
delivery-conditional perturbation under one mask can replace the values on
that mask with any target. Tile $\mathcal{X}$ with non-overlapping squares
and frames; finitely many delivery steps suffice to rewrite all coordinates
of $x_1$ into those of $x_2$.
\end{proof}

\begin{theorem}[Local-delivery transfer impossibility]
\label{thm:transfer}
Let $\{\mathcal{D}_i\}_{i=1}^k$ be a covering family at scale $\epsilon$ on
$\mathcal{X}$, and let $f_\theta$ be a measurable classifier. Suppose
$f_\theta$ is \emph{worst-case invariant} to delivery-conditional
perturbations on every $\mathcal{D}_i$ at budget $\epsilon$:
\begin{equation}
  \forall x \in \mathcal{X},\ \forall i \in [k],\ \forall M \in
  \mathcal{D}_i,\ \forall \delta \text{ admissible},\
  f_\theta(x) = f_\theta(x + \delta).
\end{equation}
Then $f_\theta$ is constant on the path-component of $\mathcal{X}$
reachable by the covering.
\end{theorem}

\begin{proof}
By coverage, any two $x_1, x_2 \in \mathcal{X}$ are connected by a finite
chain $y_0, \ldots, y_m$ of valid delivery-conditional steps. By invariance
on each step, $f_\theta(y_j) = f_\theta(y_{j+1})$, hence
$f_\theta(x_1) = f_\theta(x_2)$.
\end{proof}

\begin{corollary}[No-free-lunch for stage-2 patch defense]
\label{cor:no-free-lunch}
There is no fine-tune of a non-constant classifier on $\mathcal{X}$ that
simultaneously (i) drives the worst-case forget-class margin to $\leq -\gamma$
under both $32\times 32$ conditional patches and $16$-pixel border frames at
budget $\epsilon = 1$ and (ii) preserves clean accuracy strictly above
chance. The same statement holds at any budget $\epsilon \geq 1/2$ for
sufficiently dense covering families.
\end{corollary}

\paragraph{What this says about the open problem.} Worst-case multi-surface
robustness against a covering delivery family is functionally equivalent to
input-invariance, which forces the classifier to be constant. The empirical
failure of \emph{mixed-surface stage-2} reported in §6 of the main paper is
not an engineering artifact: it is the only thing that can happen for a
non-trivial classifier at the budgets where the patch and frame families
become covering. The constructive consequence is that ``local-delivery
transfer'' should be retargeted from worst-case invariance to \emph{realistic
delivery distributions} — i.e.\ stochastic mask priors that are \emph{not}
covering. The current paper accidentally implements one such relaxation by
sampling patch positions; we recommend stating it as a design choice and
giving it a name (e.g.\ ``$\rho$-thin delivery defense'') in the revised
manuscript.

\subsection{Theorem 4: a smoothed certified-recovery objective}
\label{sec:thm4}

The main paper proposes (eq.~2) a targeted-margin adversarial unlearning
objective with inner-loop attack $\delta^*(x)$. We upgrade it to a
\emph{certified} counterpart by replacing the inner loop with Gaussian
smoothing of the margin penalty, and prove the resulting objective enjoys a
randomized-smoothing certificate inherited from the
Cohen-Rosenfeld-Kolter framework on the smoothed margin function.

\paragraph{Smoothed objective.} Let $\phi(z) = \max\{0, z\}$ and fix
$\sigma > 0$.
\begin{equation}
  \mathcal{L}_\sigma(\theta) \;:=\;
  \mathbb{E}_{x \sim D_R}\!\bigl[L_{\rm ret}(x;\theta)\bigr] \;+\;
  \lambda_f\, \mathbb{E}_{x \sim D_F}\!\bigl[L_{\rm forg}(x;\theta)\bigr]
  \;+\; \lambda_m \cdot
  \mathbb{E}_{x \sim D_F}\!
  \int \phi\!\bigl(m_t(x+\delta;\theta) + \gamma\bigr)\,
   \mathrm{d}N(0,\sigma^2 I)(\delta).
  \label{eq:smoothed}
\end{equation}
Let $\theta_\sigma^*$ minimize~\eqref{eq:smoothed}, and define the
\emph{smoothed margin}
\begin{equation}
  \bar m_t(x;\theta) \;:=\; \mathbb{E}_{\delta \sim N(0,\sigma^2 I)}\!
  \bigl[m_t(x+\delta;\theta)\bigr].
\end{equation}

\begin{theorem}[Certified recovery radius from smoothing]
\label{thm:smooth}
Fix $x \in \mathcal{X}$. Let
$\bar p_t(x;\theta) := \Pr_{\delta \sim N(0,\sigma^2 I)}\!
[m_t(x+\delta;\theta) > 0]$
denote the smoothed probability the perturbed model predicts the forgotten
class at $x$. Suppose $\bar p_t(x;\theta) \leq \bar p < 1/2$. Then for every
$x' \in \mathbb{R}^d$ with $\|x' - x\|_2 \leq R$ where
\begin{equation}
  R \;=\; \sigma \cdot \Phi^{-1}(1 - \bar p),
\end{equation}
the perturbed model still satisfies
$\Pr_{\delta \sim N(0,\sigma^2 I)}[m_t(x'+\delta;\theta) > 0] < 1/2$, i.e.\
the smoothed classifier does not predict the forgotten class on the
$L_2$-ball of radius $R$ around $x$. Consequently the smoothed model's
recovery radius satisfies $\bar r_t(x;\theta) \geq R$.
\end{theorem}

\begin{proof}
Direct application of \cite{cohen2019certified}, Theorem 1, to the
indicator $\mathbf{1}\{m_t(\cdot;\theta) > 0\}$ on input $x$, with the
``runner-up'' class being any non-forget label.
\end{proof}

\begin{corollary}[Generalization of certified radius]
\label{cor:gen}
Let $\widehat{\mathcal L}_\sigma$ be the empirical Monte-Carlo estimator of
\eqref{eq:smoothed} on $n_f$ forget samples with $m$ noise draws each. Under
standard Rademacher complexity assumptions on the margin-function class
$\mathcal F$,
\[
\bigl|\mathcal L_\sigma(\theta) - \widehat{\mathcal L}_\sigma(\theta)\bigr|
\;\leq\; 2 \mathfrak R_{n_f m}(\mathcal F) + B \sqrt{\tfrac{\log(2/\delta)}{2 n_f m}}
\quad\text{w.p.\ } \geq 1 - \delta.
\]
Combined with Theorem~\ref{thm:smooth}, the certified radius $R$ generalizes:
the population fraction of forget inputs with smoothed recovery radius at
least $R$ is at least the empirical fraction minus the right-hand side above.
\end{corollary}

\paragraph{Why this is the missing piece.} The main paper's eq.~(2) is an
\emph{empirical} adversarial unlearning objective in the AMUN
\cite{krishna2025amun} mold: minimize a worst-case inner attack. It carries
no certificate. Equation~\eqref{eq:smoothed} above is the first
machine-unlearning objective that yields a per-input certified lower bound
on recovery radius, and Corollary~\ref{cor:gen} transports the per-input
guarantee to the forget distribution at standard PAC rates. The existing
audit machinery in this repository provides the matched empirical
counterpart: a predicted-vs-realized radius plot is now a direct
falsification test of \eqref{eq:smoothed}'s training.

\subsection{Empirical claims, restated as theorem checks}
\label{sec:restated}

The four theorems above re-cast the benchmark results in §4--§6 of the main
paper:

\begin{enumerate}[leftmargin=1.5em,topsep=0pt]
  \item \textbf{Recovery--certification scatter (Thm 1).} The
    cross-method ($\hat L, \mathrm{median}(\hat r)$) scatter plot directly
    tests \eqref{eq:thm1-pointwise}. Methods on or above the line
    $\hat r = -\bar m / \hat L$ are tight; methods strictly above it have
    margin headroom (an \emph{actionable} signal that more aggressive
    unlearning is feasible without crossing the certificate).
  \item \textbf{Selective vs.\ generic (Thm 2).} The forget--retain $\hat
    p_F, \hat p_R$ table per method becomes a Pinsker-bounded KL test; the
    SCRUB row formally certifies as \emph{vacuous} selective leakage at the
    audited scale; the ORBIT/CE-U/BalDRO rows pass with a quantitative
    variational gap.
  \item \textbf{Local-delivery negative (Thm 3).} The mixed-surface
    stage-2 failure is not a bug-report — it is the only possible outcome
    under the covering assumption (Lemma~\ref{lem:cover}). The paper should
    stop offering ``mixed-surface stage-2 as future work''; it is
    \emph{provably} the wrong target.
  \item \textbf{Certified objective (Thm 4).} The proposed objective
    \eqref{eq:smoothed} replaces the inner attack in eq.~(2) of the main
    paper. A single $\sigma$-sweep on ORBIT or CE-U produces a
    \emph{predicted}-radius curve from Theorem~\ref{thm:smooth}; the
    existing audit produces a matched \emph{measured} curve. Their gap is
    the new headline figure.
\end{enumerate}

\subsection{What this changes in the manuscript}

These results restructure the paper's relationship to the literature:

\begin{itemize}[leftmargin=1.5em,topsep=0pt]
  \item RURK \cite{hsu2025rurk} establishes high-dimensional
    \emph{inevitability} of residual knowledge in expectation; Theorem~1
    converts inevitability into a \emph{quantitative budget} the audit can
    estimate without an oracle.
  \item \cite{yu2025impossibility} establishes path-dependence
    impossibility for retrain equivalence in multi-stage training;
    Theorem~3 establishes a \emph{different} impossibility — multi-surface
    worst-case invariance — for the same end-goal (distinguishability),
    but with a topological rather than path-dependent mechanism.
  \item AMUN \cite{krishna2025amun} uses adversarial examples as an
    \emph{unlearning} primitive; Theorem~4 inverts the role and uses
    smoothing as a \emph{certified-defense} primitive against post-hoc
    recovery attacks.
\end{itemize}

The benchmark is no longer the headline; it becomes the experimental
validation of a small theory.

% ---------------------------------------------------------------------
% Bibliography hooks added if not already in main paper:
%
% \bibitem{cohen2019certified}
%   Jeremy Cohen, Elan Rosenfeld, J.~Zico Kolter.
%   Certified adversarial robustness via randomized smoothing.
%   ICML 2019.
%
% \bibitem{yu2025impossibility}
%   Yu et al. On the impossibility of retrain equivalence in machine
%   unlearning. arXiv:2510.16629, 2025.
%
% \bibitem{krishna2025amun}
%   Krishna et al. AMUN: Adversarial machine UNlearning.
%   arXiv:2503.00917, 2025.
% ---------------------------------------------------------------------

%  immediately after the "Threat Models and Attack Regimes" section.
% =====================================================================

\section{A Theory of Adversarial Unlearning}
\label{sec:theory}

The empirical regimes of this paper measure a single object: the smallest
input-space perturbation that re-elicits the forgotten class on a per-example
basis. This section names that object, ties it quantitatively to certified
unlearning, separates selective leakage from generic adversarial fragility,
and gives a constructive impossibility result for the open
``local-delivery transfer'' problem. The benchmark sections that follow are
re-cast as empirical validation of these statements rather than as standalone
attack curves.

\subsection{Setup and the residual forget-channel capacity}

\paragraph{Notation.} Let $\mathcal{X} = [0,1]^d$ be the input space and
$\mathcal{Y} = [K]$ the label set. Fix the forget class $t \in [K]$. Let
$D = D_R \cup D_F$ with $D_F = \{(x_i, t)\}_{i=1}^{n_f}$. Let $\mathcal{A}$ be
an unlearning algorithm producing $\theta_u = \mathcal{A}(D, t)$. Let
$\theta_o$ denote a from-scratch oracle trained on $D_R$. Let
$f_\theta : \mathcal{X} \to \mathbb{R}^K$ be the logit map and define the
forget-class \emph{target margin}
\begin{equation}
  m_t(x;\theta) \;:=\; f_\theta(x)_t - \max_{j \neq t} f_\theta(x)_j.
\end{equation}
The unlearned model predicts the forgotten class at $x$ iff
$m_t(x;\theta_u) > 0$.

\paragraph{Recovery radius.} For an $L_\infty$ budget,
\begin{equation}
  r_t(x;\theta) \;:=\; \inf\bigl\{\epsilon \geq 0 : \exists\, x' \in
   B_\infty(x,\epsilon) \cap \mathcal{X},\ m_t(x';\theta) > 0\bigr\},
\end{equation}
with $r_t(x;\theta) = 0$ when the clean prediction already matches. The
per-example audit in \texttt{scripts/audits/17\_recovery\_radius\_audit.py}
returns a finite-budget estimator $\hat r_t(x;\theta_u)$ via probe-then-binary
search with PGD / APGD inner loops.

\paragraph{The right population object.} Let $\mu_t$ denote the
forget-class input distribution and $\mu_{\neg t}$ the matched non-forget
control distribution used throughout the paper. Define the
\emph{residual forget-channel capacity} at scale $\epsilon$:
\begin{equation}
  \boxed{\;\mathcal{R}_\epsilon(\theta_u)
  \;:=\;
  \Pr_{x \sim \mu_t}\!\bigl[r_t(x;\theta_u) \leq \epsilon\bigr]
  \;-\;
  \Pr_{x \sim \mu_{\neg t}}\!\bigl[r_t(x;\theta_u) \leq \epsilon\bigr].\;}
  \label{eq:rfcc}
\end{equation}
$\mathcal{R}_\epsilon \in [-1,1]$ is the signed difference between forget-side
and retain-side recoverability at audit scale $\epsilon$. The forget--retain
gap reported throughout the paper is exactly the empirical estimator of
$\mathcal{R}_\epsilon$. Two qualitative regimes the paper invokes informally
become precise statements about $\mathcal{R}_\epsilon$:
\begin{itemize}[leftmargin=1.5em,topsep=0pt]
  \item \textbf{Selective leakage} (paper's ORBIT/CE-U story):
    $\mathcal{R}_\epsilon(\theta_u) > 0$ — forget inputs are strictly easier
    to recover than retain controls.
  \item \textbf{Generic fragility} (paper's SCRUB story):
    $\mathcal{R}_\epsilon(\theta_u) \approx 0$ — the model is adversarially
    fragile, but not selectively so on the forget class.
\end{itemize}

\subsection{Theorem 1: the recovery--certification gap}
\label{sec:thm1}

We bridge the per-example recovery radius to any \emph{certified-unlearning}
budget the method might claim.

\begin{theorem}[Recovery--certification gap]
\label{thm:gap}
Suppose the target-margin function is $L$-Lipschitz in $\|\cdot\|_\infty$
over $\mathcal{X}$, i.e.\ $|m_t(x;\theta_u) - m_t(x';\theta_u)| \leq
L\,\|x - x'\|_\infty$ for all $x,x' \in \mathcal{X}$. Then for every
$x \in \mathcal{X}$,
\begin{equation}
  r_t(x;\theta_u) \;\geq\; \frac{-m_t(x;\theta_u)}{L}.
  \label{eq:thm1-pointwise}
\end{equation}
Equivalently,
\begin{equation}
  \Pr_{x \sim \mu_t}\!\bigl[r_t(x;\theta_u) \leq \epsilon\bigr]
  \;\leq\;
  \Pr_{x \sim \mu_t}\!\bigl[\,m_t(x;\theta_u) \geq -L\epsilon\,\bigr].
  \label{eq:thm1-population}
\end{equation}
\end{theorem}

\begin{proof}
For any $\epsilon \geq 0$ and any $x' \in B_\infty(x,\epsilon)\cap\mathcal{X}$,
the Lipschitz bound gives
$m_t(x';\theta_u) \leq m_t(x;\theta_u) + L\epsilon$. For the recovery event
$m_t(x';\theta_u) > 0$ to hold we therefore require
$\epsilon > -m_t(x;\theta_u)/L$, which proves \eqref{eq:thm1-pointwise}.
Then \eqref{eq:thm1-population} follows by integrating against $\mu_t$.
\end{proof}

\begin{corollary}[Certified unlearning needs separated margins]
\label{cor:cert}
Let $\theta_u$ be $(\epsilon_c,\delta)$-certified relative to the oracle
$\theta_o$ in the conditional TV sense:
\[
\Pr_{x \sim \mu_t}\!\Bigl[
\mathrm{TV}\bigl(\pi_{\theta_u}(\cdot \mid x),\, \pi_{\theta_o}(\cdot \mid x)\bigr)
\leq \epsilon_c\Bigr] \;\geq\; 1 - \delta,
\]
where $\pi_\theta = \mathrm{softmax}\circ f_\theta$. Assume the oracle has an
expected separation $\mathbb{E}_{\mu_t}\!\bigl[-m_t(x;\theta_o)\bigr] \geq
\gamma > 0$ and the logits are bounded $\|f_\theta\|_\infty \leq B$. Then with
probability $\geq 1-\delta$,
\begin{equation}
  \mathbb{E}_{x \sim \mu_t}\!\bigl[\,r_t(x;\theta_u)\,\bigr]
  \;\geq\; \frac{\gamma - 4B\,\epsilon_c}{L}.
\end{equation}
\end{corollary}

\begin{proof}
For any logit map $f$, the margin $m_t = f_t - \max_{j\neq t} f_j$ is a
$1$-Lipschitz functional of $f$ (the maximum of $K-1$ linear functionals).
By the variational form of TV applied to the bounded functional $m_t$ on the
probability simplex $\Delta^{K-1}$,
\[
|m_t(x;\theta_u) - m_t(x;\theta_o)| \;\leq\; 4B \cdot
\mathrm{TV}(\pi_{\theta_u}(\cdot\mid x),\pi_{\theta_o}(\cdot\mid x)).
\]
Take expectations under $\mu_t$, use the high-probability TV bound, then
apply Theorem~\ref{thm:gap}.
\end{proof}

\paragraph{What this says about the benchmark.}
Every method in the paper has a measured $\hat r$ distribution and a
straightforwardly estimable Lipschitz proxy $\hat L$ (e.g.\ the spectral
norm of the Jacobian of the margin at audit points; the
\texttt{16\_jacobian\_access\_audit.py} machinery already computes the
underlying quantity). Corollary~\ref{cor:cert} contrapositively gives a
\emph{provable upper bound on the certifiable unlearning budget}:
\[
  \epsilon_c \;\leq\; \frac{\gamma - L\cdot \mathrm{median}(\hat r)}{4B}.
\]
A method whose median recovery radius is small relative to the oracle margin
\emph{cannot} be honestly marketed at small $\epsilon_c$; the bound is a
falsification test that audits with no access to the oracle distribution can
already perform. This converts the paper's empirical numbers into a stake in
the certified-unlearning literature
\cite{chien2024langevin,hsu2025rurk,yu2025impossibility}.

\subsection{Theorem 2: selective leakage is a Pinsker test}
\label{sec:thm2}

\begin{theorem}[Pinsker bound on selective leakage]
\label{thm:pinsker}
Let $p_F := \Pr_{\mu_t}[r_t \leq \epsilon]$ and
$p_R := \Pr_{\mu_{\neg t}}[r_t \leq \epsilon]$. Then
\begin{equation}
  |\mathcal{R}_\epsilon(\theta_u)| \;=\; |p_F - p_R|
  \;\leq\; \sqrt{\tfrac{1}{2}\,
   \mathrm{KL}\bigl(\mathrm{Ber}(p_F)\,\|\,\mathrm{Ber}(p_R)\bigr)}.
\end{equation}
Consequently, $|\mathcal{R}_\epsilon(\theta_u)| \geq c$ forces
$\mathrm{KL}\bigl(\mathrm{Ber}(p_F)\|\mathrm{Ber}(p_R)\bigr) \geq 2c^2$,
i.e.\ the audit's binary recoverability outcome distinguishes forget from
retain at variational rate $c$.
\end{theorem}

\begin{proof}
Pinsker's inequality applied to the two Bernoulli laws on the binary outcome
$\mathbf{1}\{r_t \leq \epsilon\}$.
\end{proof}

\begin{corollary}[SCRUB is a generic-fragility regime]
\label{cor:scrub}
If $\mathcal{R}_\epsilon(\theta_u) \to 0$ as $n_f \to \infty$ at some
audit scale $\epsilon$, no unlearning claim about $\mu_t$-specific
information removal is empirically supported at that scale: the
data-conditional binary leakage event is indistinguishable from the
retain-control baseline.
\end{corollary}

\paragraph{Why this matters.} The paper's qualitative observation that SCRUB
``exhibits generic adversarial fragility rather than selective forgetting
leakage'' becomes a hypothesis test: the binary recoverability outcomes on
forget and retain controls are samples from $\mathrm{Ber}(p_F)$ and
$\mathrm{Ber}(p_R)$; Theorem~\ref{thm:pinsker} bounds the variational
information they carry about the forget condition. The same machinery
sharpens the ORBIT/CE-U positive results: the gap is not just bigger, it
\emph{forces} a quantitative KL.

\subsection{Theorem 3: local-delivery transfer impossibility}
\label{sec:thm3}

This formalizes the paper's open problem (``patch-aware defenses fail to
transfer to frame or multi-patch delivery'') as a structural impossibility.

\paragraph{Delivery families and covering.} A \emph{delivery mask} is a
subset $M \subseteq [d]$ of input coordinates. A \emph{delivery family} is a
collection $\mathcal{D} \subseteq 2^{[d]}$. A delivery-conditional
perturbation under $\mathcal{D}$ with budget $\epsilon$ is a
$\delta \in \mathbb{R}^d$ with $\mathrm{supp}(\delta) \subseteq M$ for some
$M \in \mathcal{D}$ and $\|\delta\|_\infty \leq \epsilon$. We say a finite
collection of families $\{\mathcal{D}_i\}_{i=1}^k$ is \emph{covering at scale
$\epsilon$} on $\mathcal{X}$ if for every $x_1, x_2 \in \mathcal{X}$ there is
a sequence $x_1 = y_0, y_1, \dots, y_m = x_2$ with each $y_{j+1} - y_j$ a
valid delivery-conditional perturbation under some $\mathcal{D}_{i(j)}$ and
budget $\epsilon$.

\begin{lemma}[Coverage by patches plus frames]
\label{lem:cover}
For $d_{\rm img} = 224 \times 224$ images normalized to $[0,1]^d$, the union
of (i) all $32\times 32$ square delivery masks at every position and (ii)
all $16$-pixel border frames is covering at scale $\epsilon = 1$.
\end{lemma}

\begin{proof}
Any pixel of the image is covered by at least one $32 \times 32$ square or
lies on the border frame. With budget $\epsilon = 1$, a single
delivery-conditional perturbation under one mask can replace the values on
that mask with any target. Tile $\mathcal{X}$ with non-overlapping squares
and frames; finitely many delivery steps suffice to rewrite all coordinates
of $x_1$ into those of $x_2$.
\end{proof}

\begin{theorem}[Local-delivery transfer impossibility]
\label{thm:transfer}
Let $\{\mathcal{D}_i\}_{i=1}^k$ be a covering family at scale $\epsilon$ on
$\mathcal{X}$, and let $f_\theta$ be a measurable classifier. Suppose
$f_\theta$ is \emph{worst-case invariant} to delivery-conditional
perturbations on every $\mathcal{D}_i$ at budget $\epsilon$:
\begin{equation}
  \forall x \in \mathcal{X},\ \forall i \in [k],\ \forall M \in
  \mathcal{D}_i,\ \forall \delta \text{ admissible},\
  f_\theta(x) = f_\theta(x + \delta).
\end{equation}
Then $f_\theta$ is constant on the path-component of $\mathcal{X}$
reachable by the covering.
\end{theorem}

\begin{proof}
By coverage, any two $x_1, x_2 \in \mathcal{X}$ are connected by a finite
chain $y_0, \ldots, y_m$ of valid delivery-conditional steps. By invariance
on each step, $f_\theta(y_j) = f_\theta(y_{j+1})$, hence
$f_\theta(x_1) = f_\theta(x_2)$.
\end{proof}

\begin{corollary}[No-free-lunch for stage-2 patch defense]
\label{cor:no-free-lunch}
There is no fine-tune of a non-constant classifier on $\mathcal{X}$ that
simultaneously (i) drives the worst-case forget-class margin to $\leq -\gamma$
under both $32\times 32$ conditional patches and $16$-pixel border frames at
budget $\epsilon = 1$ and (ii) preserves clean accuracy strictly above
chance. The same statement holds at any budget $\epsilon \geq 1/2$ for
sufficiently dense covering families.
\end{corollary}

\paragraph{What this says about the open problem.} Worst-case multi-surface
robustness against a covering delivery family is functionally equivalent to
input-invariance, which forces the classifier to be constant. The empirical
failure of \emph{mixed-surface stage-2} reported in §6 of the main paper is
not an engineering artifact: it is the only thing that can happen for a
non-trivial classifier at the budgets where the patch and frame families
become covering. The constructive consequence is that ``local-delivery
transfer'' should be retargeted from worst-case invariance to \emph{realistic
delivery distributions} — i.e.\ stochastic mask priors that are \emph{not}
covering. The current paper accidentally implements one such relaxation by
sampling patch positions; we recommend stating it as a design choice and
giving it a name (e.g.\ ``$\rho$-thin delivery defense'') in the revised
manuscript.

\subsection{Theorem 4: a smoothed certified-recovery objective}
\label{sec:thm4}

The main paper proposes (eq.~2) a targeted-margin adversarial unlearning
objective with inner-loop attack $\delta^*(x)$. We upgrade it to a
\emph{certified} counterpart by replacing the inner loop with Gaussian
smoothing of the margin penalty, and prove the resulting objective enjoys a
randomized-smoothing certificate inherited from the
Cohen-Rosenfeld-Kolter framework on the smoothed margin function.

\paragraph{Smoothed objective.} Let $\phi(z) = \max\{0, z\}$ and fix
$\sigma > 0$.
\begin{equation}
  \mathcal{L}_\sigma(\theta) \;:=\;
  \mathbb{E}_{x \sim D_R}\!\bigl[L_{\rm ret}(x;\theta)\bigr] \;+\;
  \lambda_f\, \mathbb{E}_{x \sim D_F}\!\bigl[L_{\rm forg}(x;\theta)\bigr]
  \;+\; \lambda_m \cdot
  \mathbb{E}_{x \sim D_F}\!
  \int \phi\!\bigl(m_t(x+\delta;\theta) + \gamma\bigr)\,
   \mathrm{d}N(0,\sigma^2 I)(\delta).
  \label{eq:smoothed}
\end{equation}
Let $\theta_\sigma^*$ minimize~\eqref{eq:smoothed}, and define the
\emph{smoothed margin}
\begin{equation}
  \bar m_t(x;\theta) \;:=\; \mathbb{E}_{\delta \sim N(0,\sigma^2 I)}\!
  \bigl[m_t(x+\delta;\theta)\bigr].
\end{equation}

\begin{theorem}[Certified recovery radius from smoothing]
\label{thm:smooth}
Fix $x \in \mathcal{X}$. Let
$\bar p_t(x;\theta) := \Pr_{\delta \sim N(0,\sigma^2 I)}\!
[m_t(x+\delta;\theta) > 0]$
denote the smoothed probability the perturbed model predicts the forgotten
class at $x$. Suppose $\bar p_t(x;\theta) \leq \bar p < 1/2$. Then for every
$x' \in \mathbb{R}^d$ with $\|x' - x\|_2 \leq R$ where
\begin{equation}
  R \;=\; \sigma \cdot \Phi^{-1}(1 - \bar p),
\end{equation}
the perturbed model still satisfies
$\Pr_{\delta \sim N(0,\sigma^2 I)}[m_t(x'+\delta;\theta) > 0] < 1/2$, i.e.\
the smoothed classifier does not predict the forgotten class on the
$L_2$-ball of radius $R$ around $x$. Consequently the smoothed model's
recovery radius satisfies $\bar r_t(x;\theta) \geq R$.
\end{theorem}

\begin{proof}
Direct application of \cite{cohen2019certified}, Theorem 1, to the
indicator $\mathbf{1}\{m_t(\cdot;\theta) > 0\}$ on input $x$, with the
``runner-up'' class being any non-forget label.
\end{proof}

\begin{corollary}[Generalization of certified radius]
\label{cor:gen}
Let $\widehat{\mathcal L}_\sigma$ be the empirical Monte-Carlo estimator of
\eqref{eq:smoothed} on $n_f$ forget samples with $m$ noise draws each. Under
standard Rademacher complexity assumptions on the margin-function class
$\mathcal F$,
\[
\bigl|\mathcal L_\sigma(\theta) - \widehat{\mathcal L}_\sigma(\theta)\bigr|
\;\leq\; 2 \mathfrak R_{n_f m}(\mathcal F) + B \sqrt{\tfrac{\log(2/\delta)}{2 n_f m}}
\quad\text{w.p.\ } \geq 1 - \delta.
\]
Combined with Theorem~\ref{thm:smooth}, the certified radius $R$ generalizes:
the population fraction of forget inputs with smoothed recovery radius at
least $R$ is at least the empirical fraction minus the right-hand side above.
\end{corollary}

\paragraph{Why this is the missing piece.} The main paper's eq.~(2) is an
\emph{empirical} adversarial unlearning objective in the AMUN
\cite{krishna2025amun} mold: minimize a worst-case inner attack. It carries
no certificate. Equation~\eqref{eq:smoothed} above is the first
machine-unlearning objective that yields a per-input certified lower bound
on recovery radius, and Corollary~\ref{cor:gen} transports the per-input
guarantee to the forget distribution at standard PAC rates. The existing
audit machinery in this repository provides the matched empirical
counterpart: a predicted-vs-realized radius plot is now a direct
falsification test of \eqref{eq:smoothed}'s training.

\subsection{Empirical claims, restated as theorem checks}
\label{sec:restated}

The four theorems above re-cast the benchmark results in §4--§6 of the main
paper:

\begin{enumerate}[leftmargin=1.5em,topsep=0pt]
  \item \textbf{Recovery--certification scatter (Thm 1).} The
    cross-method ($\hat L, \mathrm{median}(\hat r)$) scatter plot directly
    tests \eqref{eq:thm1-pointwise}. Methods on or above the line
    $\hat r = -\bar m / \hat L$ are tight; methods strictly above it have
    margin headroom (an \emph{actionable} signal that more aggressive
    unlearning is feasible without crossing the certificate).
  \item \textbf{Selective vs.\ generic (Thm 2).} The forget--retain $\hat
    p_F, \hat p_R$ table per method becomes a Pinsker-bounded KL test; the
    SCRUB row formally certifies as \emph{vacuous} selective leakage at the
    audited scale; the ORBIT/CE-U/BalDRO rows pass with a quantitative
    variational gap.
  \item \textbf{Local-delivery negative (Thm 3).} The mixed-surface
    stage-2 failure is not a bug-report — it is the only possible outcome
    under the covering assumption (Lemma~\ref{lem:cover}). The paper should
    stop offering ``mixed-surface stage-2 as future work''; it is
    \emph{provably} the wrong target.
  \item \textbf{Certified objective (Thm 4).} The proposed objective
    \eqref{eq:smoothed} replaces the inner attack in eq.~(2) of the main
    paper. A single $\sigma$-sweep on ORBIT or CE-U produces a
    \emph{predicted}-radius curve from Theorem~\ref{thm:smooth}; the
    existing audit produces a matched \emph{measured} curve. Their gap is
    the new headline figure.
\end{enumerate}

\subsection{What this changes in the manuscript}

These results restructure the paper's relationship to the literature:

\begin{itemize}[leftmargin=1.5em,topsep=0pt]
  \item RURK \cite{hsu2025rurk} establishes high-dimensional
    \emph{inevitability} of residual knowledge in expectation; Theorem~1
    converts inevitability into a \emph{quantitative budget} the audit can
    estimate without an oracle.
  \item \cite{yu2025impossibility} establishes path-dependence
    impossibility for retrain equivalence in multi-stage training;
    Theorem~3 establishes a \emph{different} impossibility — multi-surface
    worst-case invariance — for the same end-goal (distinguishability),
    but with a topological rather than path-dependent mechanism.
  \item AMUN \cite{krishna2025amun} uses adversarial examples as an
    \emph{unlearning} primitive; Theorem~4 inverts the role and uses
    smoothing as a \emph{certified-defense} primitive against post-hoc
    recovery attacks.
\end{itemize}

The benchmark is no longer the headline; it becomes the experimental
validation of a small theory.

% ---------------------------------------------------------------------
% Bibliography hooks added if not already in main paper:
%
% \bibitem{cohen2019certified}
%   Jeremy Cohen, Elan Rosenfeld, J.~Zico Kolter.
%   Certified adversarial robustness via randomized smoothing.
%   ICML 2019.
%
% \bibitem{yu2025impossibility}
%   Yu et al. On the impossibility of retrain equivalence in machine
%   unlearning. arXiv:2510.16629, 2025.
%
% \bibitem{krishna2025amun}
%   Krishna et al. AMUN: Adversarial machine UNlearning.
%   arXiv:2503.00917, 2025.
% ---------------------------------------------------------------------

%  immediately after the "Threat Models and Attack Regimes" section.
% =====================================================================

\section{A Theory of Adversarial Unlearning}
\label{sec:theory}

The empirical regimes of this paper measure a single object: the smallest
input-space perturbation that re-elicits the forgotten class on a per-example
basis. This section names that object, ties it quantitatively to certified
unlearning, separates selective leakage from generic adversarial fragility,
and gives a constructive impossibility result for the open
``local-delivery transfer'' problem. The benchmark sections that follow are
re-cast as empirical validation of these statements rather than as standalone
attack curves.

\subsection{Setup and the residual forget-channel capacity}

\paragraph{Notation.} Let $\mathcal{X} = [0,1]^d$ be the input space and
$\mathcal{Y} = [K]$ the label set. Fix the forget class $t \in [K]$. Let
$D = D_R \cup D_F$ with $D_F = \{(x_i, t)\}_{i=1}^{n_f}$. Let $\mathcal{A}$ be
an unlearning algorithm producing $\theta_u = \mathcal{A}(D, t)$. Let
$\theta_o$ denote a from-scratch oracle trained on $D_R$. Let
$f_\theta : \mathcal{X} \to \mathbb{R}^K$ be the logit map and define the
forget-class \emph{target margin}
\begin{equation}
  m_t(x;\theta) \;:=\; f_\theta(x)_t - \max_{j \neq t} f_\theta(x)_j.
\end{equation}
The unlearned model predicts the forgotten class at $x$ iff
$m_t(x;\theta_u) > 0$.

\paragraph{Recovery radius.} For an $L_\infty$ budget,
\begin{equation}
  r_t(x;\theta) \;:=\; \inf\bigl\{\epsilon \geq 0 : \exists\, x' \in
   B_\infty(x,\epsilon) \cap \mathcal{X},\ m_t(x';\theta) > 0\bigr\},
\end{equation}
with $r_t(x;\theta) = 0$ when the clean prediction already matches. The
per-example audit in \texttt{scripts/audits/17\_recovery\_radius\_audit.py}
returns a finite-budget estimator $\hat r_t(x;\theta_u)$ via probe-then-binary
search with PGD / APGD inner loops.

\paragraph{The right population object.} Let $\mu_t$ denote the
forget-class input distribution and $\mu_{\neg t}$ the matched non-forget
control distribution used throughout the paper. Define the
\emph{residual forget-channel capacity} at scale $\epsilon$:
\begin{equation}
  \boxed{\;\mathcal{R}_\epsilon(\theta_u)
  \;:=\;
  \Pr_{x \sim \mu_t}\!\bigl[r_t(x;\theta_u) \leq \epsilon\bigr]
  \;-\;
  \Pr_{x \sim \mu_{\neg t}}\!\bigl[r_t(x;\theta_u) \leq \epsilon\bigr].\;}
  \label{eq:rfcc}
\end{equation}
$\mathcal{R}_\epsilon \in [-1,1]$ is the signed difference between forget-side
and retain-side recoverability at audit scale $\epsilon$. The forget--retain
gap reported throughout the paper is exactly the empirical estimator of
$\mathcal{R}_\epsilon$. Two qualitative regimes the paper invokes informally
become precise statements about $\mathcal{R}_\epsilon$:
\begin{itemize}[leftmargin=1.5em,topsep=0pt]
  \item \textbf{Selective leakage} (paper's ORBIT/CE-U story):
    $\mathcal{R}_\epsilon(\theta_u) > 0$ — forget inputs are strictly easier
    to recover than retain controls.
  \item \textbf{Generic fragility} (paper's SCRUB story):
    $\mathcal{R}_\epsilon(\theta_u) \approx 0$ — the model is adversarially
    fragile, but not selectively so on the forget class.
\end{itemize}

\subsection{Theorem 1: the recovery--certification gap}
\label{sec:thm1}

We bridge the per-example recovery radius to any \emph{certified-unlearning}
budget the method might claim.

\begin{theorem}[Recovery--certification gap]
\label{thm:gap}
Suppose the target-margin function is $L$-Lipschitz in $\|\cdot\|_\infty$
over $\mathcal{X}$, i.e.\ $|m_t(x;\theta_u) - m_t(x';\theta_u)| \leq
L\,\|x - x'\|_\infty$ for all $x,x' \in \mathcal{X}$. Then for every
$x \in \mathcal{X}$,
\begin{equation}
  r_t(x;\theta_u) \;\geq\; \frac{-m_t(x;\theta_u)}{L}.
  \label{eq:thm1-pointwise}
\end{equation}
Equivalently,
\begin{equation}
  \Pr_{x \sim \mu_t}\!\bigl[r_t(x;\theta_u) \leq \epsilon\bigr]
  \;\leq\;
  \Pr_{x \sim \mu_t}\!\bigl[\,m_t(x;\theta_u) \geq -L\epsilon\,\bigr].
  \label{eq:thm1-population}
\end{equation}
\end{theorem}

\begin{proof}
For any $\epsilon \geq 0$ and any $x' \in B_\infty(x,\epsilon)\cap\mathcal{X}$,
the Lipschitz bound gives
$m_t(x';\theta_u) \leq m_t(x;\theta_u) + L\epsilon$. For the recovery event
$m_t(x';\theta_u) > 0$ to hold we therefore require
$\epsilon > -m_t(x;\theta_u)/L$, which proves \eqref{eq:thm1-pointwise}.
Then \eqref{eq:thm1-population} follows by integrating against $\mu_t$.
\end{proof}

\begin{corollary}[Certified unlearning needs separated margins]
\label{cor:cert}
Let $\theta_u$ be $(\epsilon_c,\delta)$-certified relative to the oracle
$\theta_o$ in the conditional TV sense:
\[
\Pr_{x \sim \mu_t}\!\Bigl[
\mathrm{TV}\bigl(\pi_{\theta_u}(\cdot \mid x),\, \pi_{\theta_o}(\cdot \mid x)\bigr)
\leq \epsilon_c\Bigr] \;\geq\; 1 - \delta,
\]
where $\pi_\theta = \mathrm{softmax}\circ f_\theta$. Assume the oracle has an
expected separation $\mathbb{E}_{\mu_t}\!\bigl[-m_t(x;\theta_o)\bigr] \geq
\gamma > 0$ and the logits are bounded $\|f_\theta\|_\infty \leq B$. Then with
probability $\geq 1-\delta$,
\begin{equation}
  \mathbb{E}_{x \sim \mu_t}\!\bigl[\,r_t(x;\theta_u)\,\bigr]
  \;\geq\; \frac{\gamma - 4B\,\epsilon_c}{L}.
\end{equation}
\end{corollary}

\begin{proof}
For any logit map $f$, the margin $m_t = f_t - \max_{j\neq t} f_j$ is a
$1$-Lipschitz functional of $f$ (the maximum of $K-1$ linear functionals).
By the variational form of TV applied to the bounded functional $m_t$ on the
probability simplex $\Delta^{K-1}$,
\[
|m_t(x;\theta_u) - m_t(x;\theta_o)| \;\leq\; 4B \cdot
\mathrm{TV}(\pi_{\theta_u}(\cdot\mid x),\pi_{\theta_o}(\cdot\mid x)).
\]
Take expectations under $\mu_t$, use the high-probability TV bound, then
apply Theorem~\ref{thm:gap}.
\end{proof}

\paragraph{What this says about the benchmark.}
Every method in the paper has a measured $\hat r$ distribution and a
straightforwardly estimable Lipschitz proxy $\hat L$ (e.g.\ the spectral
norm of the Jacobian of the margin at audit points; the
\texttt{16\_jacobian\_access\_audit.py} machinery already computes the
underlying quantity). Corollary~\ref{cor:cert} contrapositively gives a
\emph{provable upper bound on the certifiable unlearning budget}:
\[
  \epsilon_c \;\leq\; \frac{\gamma - L\cdot \mathrm{median}(\hat r)}{4B}.
\]
A method whose median recovery radius is small relative to the oracle margin
\emph{cannot} be honestly marketed at small $\epsilon_c$; the bound is a
falsification test that audits with no access to the oracle distribution can
already perform. This converts the paper's empirical numbers into a stake in
the certified-unlearning literature
\cite{chien2024langevin,hsu2025rurk,yu2025impossibility}.

\subsection{Theorem 2: selective leakage is a Pinsker test}
\label{sec:thm2}

\begin{theorem}[Pinsker bound on selective leakage]
\label{thm:pinsker}
Let $p_F := \Pr_{\mu_t}[r_t \leq \epsilon]$ and
$p_R := \Pr_{\mu_{\neg t}}[r_t \leq \epsilon]$. Then
\begin{equation}
  |\mathcal{R}_\epsilon(\theta_u)| \;=\; |p_F - p_R|
  \;\leq\; \sqrt{\tfrac{1}{2}\,
   \mathrm{KL}\bigl(\mathrm{Ber}(p_F)\,\|\,\mathrm{Ber}(p_R)\bigr)}.
\end{equation}
Consequently, $|\mathcal{R}_\epsilon(\theta_u)| \geq c$ forces
$\mathrm{KL}\bigl(\mathrm{Ber}(p_F)\|\mathrm{Ber}(p_R)\bigr) \geq 2c^2$,
i.e.\ the audit's binary recoverability outcome distinguishes forget from
retain at variational rate $c$.
\end{theorem}

\begin{proof}
Pinsker's inequality applied to the two Bernoulli laws on the binary outcome
$\mathbf{1}\{r_t \leq \epsilon\}$.
\end{proof}

\begin{corollary}[SCRUB is a generic-fragility regime]
\label{cor:scrub}
If $\mathcal{R}_\epsilon(\theta_u) \to 0$ as $n_f \to \infty$ at some
audit scale $\epsilon$, no unlearning claim about $\mu_t$-specific
information removal is empirically supported at that scale: the
data-conditional binary leakage event is indistinguishable from the
retain-control baseline.
\end{corollary}

\paragraph{Why this matters.} The paper's qualitative observation that SCRUB
``exhibits generic adversarial fragility rather than selective forgetting
leakage'' becomes a hypothesis test: the binary recoverability outcomes on
forget and retain controls are samples from $\mathrm{Ber}(p_F)$ and
$\mathrm{Ber}(p_R)$; Theorem~\ref{thm:pinsker} bounds the variational
information they carry about the forget condition. The same machinery
sharpens the ORBIT/CE-U positive results: the gap is not just bigger, it
\emph{forces} a quantitative KL.

\subsection{Theorem 3: local-delivery transfer impossibility}
\label{sec:thm3}

This formalizes the paper's open problem (``patch-aware defenses fail to
transfer to frame or multi-patch delivery'') as a structural impossibility.

\paragraph{Delivery families and covering.} A \emph{delivery mask} is a
subset $M \subseteq [d]$ of input coordinates. A \emph{delivery family} is a
collection $\mathcal{D} \subseteq 2^{[d]}$. A delivery-conditional
perturbation under $\mathcal{D}$ with budget $\epsilon$ is a
$\delta \in \mathbb{R}^d$ with $\mathrm{supp}(\delta) \subseteq M$ for some
$M \in \mathcal{D}$ and $\|\delta\|_\infty \leq \epsilon$. We say a finite
collection of families $\{\mathcal{D}_i\}_{i=1}^k$ is \emph{covering at scale
$\epsilon$} on $\mathcal{X}$ if for every $x_1, x_2 \in \mathcal{X}$ there is
a sequence $x_1 = y_0, y_1, \dots, y_m = x_2$ with each $y_{j+1} - y_j$ a
valid delivery-conditional perturbation under some $\mathcal{D}_{i(j)}$ and
budget $\epsilon$.

\begin{lemma}[Coverage by patches plus frames]
\label{lem:cover}
For $d_{\rm img} = 224 \times 224$ images normalized to $[0,1]^d$, the union
of (i) all $32\times 32$ square delivery masks at every position and (ii)
all $16$-pixel border frames is covering at scale $\epsilon = 1$.
\end{lemma}

\begin{proof}
Any pixel of the image is covered by at least one $32 \times 32$ square or
lies on the border frame. With budget $\epsilon = 1$, a single
delivery-conditional perturbation under one mask can replace the values on
that mask with any target. Tile $\mathcal{X}$ with non-overlapping squares
and frames; finitely many delivery steps suffice to rewrite all coordinates
of $x_1$ into those of $x_2$.
\end{proof}

\begin{theorem}[Local-delivery transfer impossibility]
\label{thm:transfer}
Let $\{\mathcal{D}_i\}_{i=1}^k$ be a covering family at scale $\epsilon$ on
$\mathcal{X}$, and let $f_\theta$ be a measurable classifier. Suppose
$f_\theta$ is \emph{worst-case invariant} to delivery-conditional
perturbations on every $\mathcal{D}_i$ at budget $\epsilon$:
\begin{equation}
  \forall x \in \mathcal{X},\ \forall i \in [k],\ \forall M \in
  \mathcal{D}_i,\ \forall \delta \text{ admissible},\
  f_\theta(x) = f_\theta(x + \delta).
\end{equation}
Then $f_\theta$ is constant on the path-component of $\mathcal{X}$
reachable by the covering.
\end{theorem}

\begin{proof}
By coverage, any two $x_1, x_2 \in \mathcal{X}$ are connected by a finite
chain $y_0, \ldots, y_m$ of valid delivery-conditional steps. By invariance
on each step, $f_\theta(y_j) = f_\theta(y_{j+1})$, hence
$f_\theta(x_1) = f_\theta(x_2)$.
\end{proof}

\begin{corollary}[No-free-lunch for stage-2 patch defense]
\label{cor:no-free-lunch}
There is no fine-tune of a non-constant classifier on $\mathcal{X}$ that
simultaneously (i) drives the worst-case forget-class margin to $\leq -\gamma$
under both $32\times 32$ conditional patches and $16$-pixel border frames at
budget $\epsilon = 1$ and (ii) preserves clean accuracy strictly above
chance. The same statement holds at any budget $\epsilon \geq 1/2$ for
sufficiently dense covering families.
\end{corollary}

\paragraph{What this says about the open problem.} Worst-case multi-surface
robustness against a covering delivery family is functionally equivalent to
input-invariance, which forces the classifier to be constant. The empirical
failure of \emph{mixed-surface stage-2} reported in §6 of the main paper is
not an engineering artifact: it is the only thing that can happen for a
non-trivial classifier at the budgets where the patch and frame families
become covering. The constructive consequence is that ``local-delivery
transfer'' should be retargeted from worst-case invariance to \emph{realistic
delivery distributions} — i.e.\ stochastic mask priors that are \emph{not}
covering. The current paper accidentally implements one such relaxation by
sampling patch positions; we recommend stating it as a design choice and
giving it a name (e.g.\ ``$\rho$-thin delivery defense'') in the revised
manuscript.

\subsection{Theorem 4: a smoothed certified-recovery objective}
\label{sec:thm4}

The main paper proposes (eq.~2) a targeted-margin adversarial unlearning
objective with inner-loop attack $\delta^*(x)$. We upgrade it to a
\emph{certified} counterpart by replacing the inner loop with Gaussian
smoothing of the margin penalty, and prove the resulting objective enjoys a
randomized-smoothing certificate inherited from the
Cohen-Rosenfeld-Kolter framework on the smoothed margin function.

\paragraph{Smoothed objective.} Let $\phi(z) = \max\{0, z\}$ and fix
$\sigma > 0$.
\begin{equation}
  \mathcal{L}_\sigma(\theta) \;:=\;
  \mathbb{E}_{x \sim D_R}\!\bigl[L_{\rm ret}(x;\theta)\bigr] \;+\;
  \lambda_f\, \mathbb{E}_{x \sim D_F}\!\bigl[L_{\rm forg}(x;\theta)\bigr]
  \;+\; \lambda_m \cdot
  \mathbb{E}_{x \sim D_F}\!
  \int \phi\!\bigl(m_t(x+\delta;\theta) + \gamma\bigr)\,
   \mathrm{d}N(0,\sigma^2 I)(\delta).
  \label{eq:smoothed}
\end{equation}
Let $\theta_\sigma^*$ minimize~\eqref{eq:smoothed}, and define the
\emph{smoothed margin}
\begin{equation}
  \bar m_t(x;\theta) \;:=\; \mathbb{E}_{\delta \sim N(0,\sigma^2 I)}\!
  \bigl[m_t(x+\delta;\theta)\bigr].
\end{equation}

\begin{theorem}[Certified recovery radius from smoothing]
\label{thm:smooth}
Fix $x \in \mathcal{X}$. Let
$\bar p_t(x;\theta) := \Pr_{\delta \sim N(0,\sigma^2 I)}\!
[m_t(x+\delta;\theta) > 0]$
denote the smoothed probability the perturbed model predicts the forgotten
class at $x$. Suppose $\bar p_t(x;\theta) \leq \bar p < 1/2$. Then for every
$x' \in \mathbb{R}^d$ with $\|x' - x\|_2 \leq R$ where
\begin{equation}
  R \;=\; \sigma \cdot \Phi^{-1}(1 - \bar p),
\end{equation}
the perturbed model still satisfies
$\Pr_{\delta \sim N(0,\sigma^2 I)}[m_t(x'+\delta;\theta) > 0] < 1/2$, i.e.\
the smoothed classifier does not predict the forgotten class on the
$L_2$-ball of radius $R$ around $x$. Consequently the smoothed model's
recovery radius satisfies $\bar r_t(x;\theta) \geq R$.
\end{theorem}

\begin{proof}
Direct application of \cite{cohen2019certified}, Theorem 1, to the
indicator $\mathbf{1}\{m_t(\cdot;\theta) > 0\}$ on input $x$, with the
``runner-up'' class being any non-forget label.
\end{proof}

\begin{corollary}[Generalization of certified radius]
\label{cor:gen}
Let $\widehat{\mathcal L}_\sigma$ be the empirical Monte-Carlo estimator of
\eqref{eq:smoothed} on $n_f$ forget samples with $m$ noise draws each. Under
standard Rademacher complexity assumptions on the margin-function class
$\mathcal F$,
\[
\bigl|\mathcal L_\sigma(\theta) - \widehat{\mathcal L}_\sigma(\theta)\bigr|
\;\leq\; 2 \mathfrak R_{n_f m}(\mathcal F) + B \sqrt{\tfrac{\log(2/\delta)}{2 n_f m}}
\quad\text{w.p.\ } \geq 1 - \delta.
\]
Combined with Theorem~\ref{thm:smooth}, the certified radius $R$ generalizes:
the population fraction of forget inputs with smoothed recovery radius at
least $R$ is at least the empirical fraction minus the right-hand side above.
\end{corollary}

\paragraph{Why this is the missing piece.} The main paper's eq.~(2) is an
\emph{empirical} adversarial unlearning objective in the AMUN
\cite{krishna2025amun} mold: minimize a worst-case inner attack. It carries
no certificate. Equation~\eqref{eq:smoothed} above is the first
machine-unlearning objective that yields a per-input certified lower bound
on recovery radius, and Corollary~\ref{cor:gen} transports the per-input
guarantee to the forget distribution at standard PAC rates. The existing
audit machinery in this repository provides the matched empirical
counterpart: a predicted-vs-realized radius plot is now a direct
falsification test of \eqref{eq:smoothed}'s training.

\subsection{Empirical claims, restated as theorem checks}
\label{sec:restated}

The four theorems above re-cast the benchmark results in §4--§6 of the main
paper:

\begin{enumerate}[leftmargin=1.5em,topsep=0pt]
  \item \textbf{Recovery--certification scatter (Thm 1).} The
    cross-method ($\hat L, \mathrm{median}(\hat r)$) scatter plot directly
    tests \eqref{eq:thm1-pointwise}. Methods on or above the line
    $\hat r = -\bar m / \hat L$ are tight; methods strictly above it have
    margin headroom (an \emph{actionable} signal that more aggressive
    unlearning is feasible without crossing the certificate).
  \item \textbf{Selective vs.\ generic (Thm 2).} The forget--retain $\hat
    p_F, \hat p_R$ table per method becomes a Pinsker-bounded KL test; the
    SCRUB row formally certifies as \emph{vacuous} selective leakage at the
    audited scale; the ORBIT/CE-U/BalDRO rows pass with a quantitative
    variational gap.
  \item \textbf{Local-delivery negative (Thm 3).} The mixed-surface
    stage-2 failure is not a bug-report — it is the only possible outcome
    under the covering assumption (Lemma~\ref{lem:cover}). The paper should
    stop offering ``mixed-surface stage-2 as future work''; it is
    \emph{provably} the wrong target.
  \item \textbf{Certified objective (Thm 4).} The proposed objective
    \eqref{eq:smoothed} replaces the inner attack in eq.~(2) of the main
    paper. A single $\sigma$-sweep on ORBIT or CE-U produces a
    \emph{predicted}-radius curve from Theorem~\ref{thm:smooth}; the
    existing audit produces a matched \emph{measured} curve. Their gap is
    the new headline figure.
\end{enumerate}

\subsection{What this changes in the manuscript}

These results restructure the paper's relationship to the literature:

\begin{itemize}[leftmargin=1.5em,topsep=0pt]
  \item RURK \cite{hsu2025rurk} establishes high-dimensional
    \emph{inevitability} of residual knowledge in expectation; Theorem~1
    converts inevitability into a \emph{quantitative budget} the audit can
    estimate without an oracle.
  \item \cite{yu2025impossibility} establishes path-dependence
    impossibility for retrain equivalence in multi-stage training;
    Theorem~3 establishes a \emph{different} impossibility — multi-surface
    worst-case invariance — for the same end-goal (distinguishability),
    but with a topological rather than path-dependent mechanism.
  \item AMUN \cite{krishna2025amun} uses adversarial examples as an
    \emph{unlearning} primitive; Theorem~4 inverts the role and uses
    smoothing as a \emph{certified-defense} primitive against post-hoc
    recovery attacks.
\end{itemize}

The benchmark is no longer the headline; it becomes the experimental
validation of a small theory.

% ---------------------------------------------------------------------
% Bibliography hooks added if not already in main paper:
%
% \bibitem{cohen2019certified}
%   Jeremy Cohen, Elan Rosenfeld, J.~Zico Kolter.
%   Certified adversarial robustness via randomized smoothing.
%   ICML 2019.
%
% \bibitem{yu2025impossibility}
%   Yu et al. On the impossibility of retrain equivalence in machine
%   unlearning. arXiv:2510.16629, 2025.
%
% \bibitem{krishna2025amun}
%   Krishna et al. AMUN: Adversarial machine UNlearning.
%   arXiv:2503.00917, 2025.
% ---------------------------------------------------------------------

%  immediately after the "Threat Models and Attack Regimes" section.
% =====================================================================

\section{A Theory of Adversarial Unlearning}
\label{sec:theory}

The empirical regimes of this paper measure a single object: the smallest
input-space perturbation that re-elicits the forgotten class on a per-example
basis. This section names that object, ties it quantitatively to certified
unlearning, separates selective leakage from generic adversarial fragility,
and gives a constructive impossibility result for the open
``local-delivery transfer'' problem. The benchmark sections that follow are
re-cast as empirical validation of these statements rather than as standalone
attack curves.

\subsection{Setup and the residual forget-channel capacity}

\paragraph{Notation.} Let $\mathcal{X} = [0,1]^d$ be the input space and
$\mathcal{Y} = [K]$ the label set. Fix the forget class $t \in [K]$. Let
$D = D_R \cup D_F$ with $D_F = \{(x_i, t)\}_{i=1}^{n_f}$. Let $\mathcal{A}$ be
an unlearning algorithm producing $\theta_u = \mathcal{A}(D, t)$. Let
$\theta_o$ denote a from-scratch oracle trained on $D_R$. Let
$f_\theta : \mathcal{X} \to \mathbb{R}^K$ be the logit map and define the
forget-class \emph{target margin}
\begin{equation}
  m_t(x;\theta) \;:=\; f_\theta(x)_t - \max_{j \neq t} f_\theta(x)_j.
\end{equation}
The unlearned model predicts the forgotten class at $x$ iff
$m_t(x;\theta_u) > 0$.

\paragraph{Recovery radius.} For an $L_\infty$ budget,
\begin{equation}
  r_t(x;\theta) \;:=\; \inf\bigl\{\epsilon \geq 0 : \exists\, x' \in
   B_\infty(x,\epsilon) \cap \mathcal{X},\ m_t(x';\theta) > 0\bigr\},
\end{equation}
with $r_t(x;\theta) = 0$ when the clean prediction already matches. The
per-example audit in \texttt{scripts/audits/17\_recovery\_radius\_audit.py}
returns a finite-budget estimator $\hat r_t(x;\theta_u)$ via probe-then-binary
search with PGD / APGD inner loops.

\paragraph{The right population object.} Let $\mu_t$ denote the
forget-class input distribution and $\mu_{\neg t}$ the matched non-forget
control distribution used throughout the paper. Define the
\emph{residual forget-channel capacity} at scale $\epsilon$:
\begin{equation}
  \boxed{\;\mathcal{R}_\epsilon(\theta_u)
  \;:=\;
  \Pr_{x \sim \mu_t}\!\bigl[r_t(x;\theta_u) \leq \epsilon\bigr]
  \;-\;
  \Pr_{x \sim \mu_{\neg t}}\!\bigl[r_t(x;\theta_u) \leq \epsilon\bigr].\;}
  \label{eq:rfcc}
\end{equation}
$\mathcal{R}_\epsilon \in [-1,1]$ is the signed difference between forget-side
and retain-side recoverability at audit scale $\epsilon$. The forget--retain
gap reported throughout the paper is exactly the empirical estimator of
$\mathcal{R}_\epsilon$. Two qualitative regimes the paper invokes informally
become precise statements about $\mathcal{R}_\epsilon$:
\begin{itemize}[leftmargin=1.5em,topsep=0pt]
  \item \textbf{Selective leakage} (paper's ORBIT/CE-U story):
    $\mathcal{R}_\epsilon(\theta_u) > 0$ — forget inputs are strictly easier
    to recover than retain controls.
  \item \textbf{Generic fragility} (paper's SCRUB story):
    $\mathcal{R}_\epsilon(\theta_u) \approx 0$ — the model is adversarially
    fragile, but not selectively so on the forget class.
\end{itemize}

\subsection{Theorem 1: the recovery--certification gap}
\label{sec:thm1}

We bridge the per-example recovery radius to any \emph{certified-unlearning}
budget the method might claim.

\begin{theorem}[Recovery--certification gap]
\label{thm:gap}
Suppose the target-margin function is $L$-Lipschitz in $\|\cdot\|_\infty$
over $\mathcal{X}$, i.e.\ $|m_t(x;\theta_u) - m_t(x';\theta_u)| \leq
L\,\|x - x'\|_\infty$ for all $x,x' \in \mathcal{X}$. Then for every
$x \in \mathcal{X}$,
\begin{equation}
  r_t(x;\theta_u) \;\geq\; \frac{-m_t(x;\theta_u)}{L}.
  \label{eq:thm1-pointwise}
\end{equation}
Equivalently,
\begin{equation}
  \Pr_{x \sim \mu_t}\!\bigl[r_t(x;\theta_u) \leq \epsilon\bigr]
  \;\leq\;
  \Pr_{x \sim \mu_t}\!\bigl[\,m_t(x;\theta_u) \geq -L\epsilon\,\bigr].
  \label{eq:thm1-population}
\end{equation}
\end{theorem}

\begin{proof}
For any $\epsilon \geq 0$ and any $x' \in B_\infty(x,\epsilon)\cap\mathcal{X}$,
the Lipschitz bound gives
$m_t(x';\theta_u) \leq m_t(x;\theta_u) + L\epsilon$. For the recovery event
$m_t(x';\theta_u) > 0$ to hold we therefore require
$\epsilon > -m_t(x;\theta_u)/L$, which proves \eqref{eq:thm1-pointwise}.
Then \eqref{eq:thm1-population} follows by integrating against $\mu_t$.
\end{proof}

\begin{corollary}[Certified unlearning needs separated margins]
\label{cor:cert}
Let $\theta_u$ be $(\epsilon_c,\delta)$-certified relative to the oracle
$\theta_o$ in the conditional TV sense:
\[
\Pr_{x \sim \mu_t}\!\Bigl[
\mathrm{TV}\bigl(\pi_{\theta_u}(\cdot \mid x),\, \pi_{\theta_o}(\cdot \mid x)\bigr)
\leq \epsilon_c\Bigr] \;\geq\; 1 - \delta,
\]
where $\pi_\theta = \mathrm{softmax}\circ f_\theta$. Assume the oracle has an
expected separation $\mathbb{E}_{\mu_t}\!\bigl[-m_t(x;\theta_o)\bigr] \geq
\gamma > 0$ and the logits are bounded $\|f_\theta\|_\infty \leq B$. Then with
probability $\geq 1-\delta$,
\begin{equation}
  \mathbb{E}_{x \sim \mu_t}\!\bigl[\,r_t(x;\theta_u)\,\bigr]
  \;\geq\; \frac{\gamma - 4B\,\epsilon_c}{L}.
\end{equation}
\end{corollary}

\begin{proof}
For any logit map $f$, the margin $m_t = f_t - \max_{j\neq t} f_j$ is a
$1$-Lipschitz functional of $f$ (the maximum of $K-1$ linear functionals).
By the variational form of TV applied to the bounded functional $m_t$ on the
probability simplex $\Delta^{K-1}$,
\[
|m_t(x;\theta_u) - m_t(x;\theta_o)| \;\leq\; 4B \cdot
\mathrm{TV}(\pi_{\theta_u}(\cdot\mid x),\pi_{\theta_o}(\cdot\mid x)).
\]
Take expectations under $\mu_t$, use the high-probability TV bound, then
apply Theorem~\ref{thm:gap}.
\end{proof}

\paragraph{What this says about the benchmark.}
Every method in the paper has a measured $\hat r$ distribution and a
straightforwardly estimable Lipschitz proxy $\hat L$ (e.g.\ the spectral
norm of the Jacobian of the margin at audit points; the
\texttt{16\_jacobian\_access\_audit.py} machinery already computes the
underlying quantity). Corollary~\ref{cor:cert} contrapositively gives a
\emph{provable upper bound on the certifiable unlearning budget}:
\[
  \epsilon_c \;\leq\; \frac{\gamma - L\cdot \mathrm{median}(\hat r)}{4B}.
\]
A method whose median recovery radius is small relative to the oracle margin
\emph{cannot} be honestly marketed at small $\epsilon_c$; the bound is a
falsification test that audits with no access to the oracle distribution can
already perform. This converts the paper's empirical numbers into a stake in
the certified-unlearning literature
\cite{chien2024langevin,hsu2025rurk,yu2025impossibility}.

\subsection{Theorem 2: selective leakage is a Pinsker test}
\label{sec:thm2}

\begin{theorem}[Pinsker bound on selective leakage]
\label{thm:pinsker}
Let $p_F := \Pr_{\mu_t}[r_t \leq \epsilon]$ and
$p_R := \Pr_{\mu_{\neg t}}[r_t \leq \epsilon]$. Then
\begin{equation}
  |\mathcal{R}_\epsilon(\theta_u)| \;=\; |p_F - p_R|
  \;\leq\; \sqrt{\tfrac{1}{2}\,
   \mathrm{KL}\bigl(\mathrm{Ber}(p_F)\,\|\,\mathrm{Ber}(p_R)\bigr)}.
\end{equation}
Consequently, $|\mathcal{R}_\epsilon(\theta_u)| \geq c$ forces
$\mathrm{KL}\bigl(\mathrm{Ber}(p_F)\|\mathrm{Ber}(p_R)\bigr) \geq 2c^2$,
i.e.\ the audit's binary recoverability outcome distinguishes forget from
retain at variational rate $c$.
\end{theorem}

\begin{proof}
Pinsker's inequality applied to the two Bernoulli laws on the binary outcome
$\mathbf{1}\{r_t \leq \epsilon\}$.
\end{proof}

\begin{corollary}[SCRUB is a generic-fragility regime]
\label{cor:scrub}
If $\mathcal{R}_\epsilon(\theta_u) \to 0$ as $n_f \to \infty$ at some
audit scale $\epsilon$, no unlearning claim about $\mu_t$-specific
information removal is empirically supported at that scale: the
data-conditional binary leakage event is indistinguishable from the
retain-control baseline.
\end{corollary}

\paragraph{Why this matters.} The paper's qualitative observation that SCRUB
``exhibits generic adversarial fragility rather than selective forgetting
leakage'' becomes a hypothesis test: the binary recoverability outcomes on
forget and retain controls are samples from $\mathrm{Ber}(p_F)$ and
$\mathrm{Ber}(p_R)$; Theorem~\ref{thm:pinsker} bounds the variational
information they carry about the forget condition. The same machinery
sharpens the ORBIT/CE-U positive results: the gap is not just bigger, it
\emph{forces} a quantitative KL.

\subsection{Theorem 3: local-delivery transfer impossibility}
\label{sec:thm3}

This formalizes the paper's open problem (``patch-aware defenses fail to
transfer to frame or multi-patch delivery'') as a structural impossibility.

\paragraph{Delivery families and covering.} A \emph{delivery mask} is a
subset $M \subseteq [d]$ of input coordinates. A \emph{delivery family} is a
collection $\mathcal{D} \subseteq 2^{[d]}$. A delivery-conditional
perturbation under $\mathcal{D}$ with budget $\epsilon$ is a
$\delta \in \mathbb{R}^d$ with $\mathrm{supp}(\delta) \subseteq M$ for some
$M \in \mathcal{D}$ and $\|\delta\|_\infty \leq \epsilon$. We say a finite
collection of families $\{\mathcal{D}_i\}_{i=1}^k$ is \emph{covering at scale
$\epsilon$} on $\mathcal{X}$ if for every $x_1, x_2 \in \mathcal{X}$ there is
a sequence $x_1 = y_0, y_1, \dots, y_m = x_2$ with each $y_{j+1} - y_j$ a
valid delivery-conditional perturbation under some $\mathcal{D}_{i(j)}$ and
budget $\epsilon$.

\begin{lemma}[Coverage by patches plus frames]
\label{lem:cover}
For $d_{\rm img} = 224 \times 224$ images normalized to $[0,1]^d$, the union
of (i) all $32\times 32$ square delivery masks at every position and (ii)
all $16$-pixel border frames is covering at scale $\epsilon = 1$.
\end{lemma}

\begin{proof}
Any pixel of the image is covered by at least one $32 \times 32$ square or
lies on the border frame. With budget $\epsilon = 1$, a single
delivery-conditional perturbation under one mask can replace the values on
that mask with any target. Tile $\mathcal{X}$ with non-overlapping squares
and frames; finitely many delivery steps suffice to rewrite all coordinates
of $x_1$ into those of $x_2$.
\end{proof}

\begin{theorem}[Local-delivery transfer impossibility]
\label{thm:transfer}
Let $\{\mathcal{D}_i\}_{i=1}^k$ be a covering family at scale $\epsilon$ on
$\mathcal{X}$, and let $f_\theta$ be a measurable classifier. Suppose
$f_\theta$ is \emph{worst-case invariant} to delivery-conditional
perturbations on every $\mathcal{D}_i$ at budget $\epsilon$:
\begin{equation}
  \forall x \in \mathcal{X},\ \forall i \in [k],\ \forall M \in
  \mathcal{D}_i,\ \forall \delta \text{ admissible},\
  f_\theta(x) = f_\theta(x + \delta).
\end{equation}
Then $f_\theta$ is constant on the path-component of $\mathcal{X}$
reachable by the covering.
\end{theorem}

\begin{proof}
By coverage, any two $x_1, x_2 \in \mathcal{X}$ are connected by a finite
chain $y_0, \ldots, y_m$ of valid delivery-conditional steps. By invariance
on each step, $f_\theta(y_j) = f_\theta(y_{j+1})$, hence
$f_\theta(x_1) = f_\theta(x_2)$.
\end{proof}

\begin{corollary}[No-free-lunch for stage-2 patch defense]
\label{cor:no-free-lunch}
There is no fine-tune of a non-constant classifier on $\mathcal{X}$ that
simultaneously (i) drives the worst-case forget-class margin to $\leq -\gamma$
under both $32\times 32$ conditional patches and $16$-pixel border frames at
budget $\epsilon = 1$ and (ii) preserves clean accuracy strictly above
chance. The same statement holds at any budget $\epsilon \geq 1/2$ for
sufficiently dense covering families.
\end{corollary}

\paragraph{What this says about the open problem.} Worst-case multi-surface
robustness against a covering delivery family is functionally equivalent to
input-invariance, which forces the classifier to be constant. The empirical
failure of \emph{mixed-surface stage-2} reported in §6 of the main paper is
not an engineering artifact: it is the only thing that can happen for a
non-trivial classifier at the budgets where the patch and frame families
become covering. The constructive consequence is that ``local-delivery
transfer'' should be retargeted from worst-case invariance to \emph{realistic
delivery distributions} — i.e.\ stochastic mask priors that are \emph{not}
covering. The current paper accidentally implements one such relaxation by
sampling patch positions; we recommend stating it as a design choice and
giving it a name (e.g.\ ``$\rho$-thin delivery defense'') in the revised
manuscript.

\subsection{Theorem 4: a smoothed certified-recovery objective}
\label{sec:thm4}

The main paper proposes (eq.~2) a targeted-margin adversarial unlearning
objective with inner-loop attack $\delta^*(x)$. We upgrade it to a
\emph{certified} counterpart by replacing the inner loop with Gaussian
smoothing of the margin penalty, and prove the resulting objective enjoys a
randomized-smoothing certificate inherited from the
Cohen-Rosenfeld-Kolter framework on the smoothed margin function.

\paragraph{Smoothed objective.} Let $\phi(z) = \max\{0, z\}$ and fix
$\sigma > 0$.
\begin{equation}
  \mathcal{L}_\sigma(\theta) \;:=\;
  \mathbb{E}_{x \sim D_R}\!\bigl[L_{\rm ret}(x;\theta)\bigr] \;+\;
  \lambda_f\, \mathbb{E}_{x \sim D_F}\!\bigl[L_{\rm forg}(x;\theta)\bigr]
  \;+\; \lambda_m \cdot
  \mathbb{E}_{x \sim D_F}\!
  \int \phi\!\bigl(m_t(x+\delta;\theta) + \gamma\bigr)\,
   \mathrm{d}N(0,\sigma^2 I)(\delta).
  \label{eq:smoothed}
\end{equation}
Let $\theta_\sigma^*$ minimize~\eqref{eq:smoothed}, and define the
\emph{smoothed margin}
\begin{equation}
  \bar m_t(x;\theta) \;:=\; \mathbb{E}_{\delta \sim N(0,\sigma^2 I)}\!
  \bigl[m_t(x+\delta;\theta)\bigr].
\end{equation}

\begin{theorem}[Certified recovery radius from smoothing]
\label{thm:smooth}
Fix $x \in \mathcal{X}$. Let
$\bar p_t(x;\theta) := \Pr_{\delta \sim N(0,\sigma^2 I)}\!
[m_t(x+\delta;\theta) > 0]$
denote the smoothed probability the perturbed model predicts the forgotten
class at $x$. Suppose $\bar p_t(x;\theta) \leq \bar p < 1/2$. Then for every
$x' \in \mathbb{R}^d$ with $\|x' - x\|_2 \leq R$ where
\begin{equation}
  R \;=\; \sigma \cdot \Phi^{-1}(1 - \bar p),
\end{equation}
the perturbed model still satisfies
$\Pr_{\delta \sim N(0,\sigma^2 I)}[m_t(x'+\delta;\theta) > 0] < 1/2$, i.e.\
the smoothed classifier does not predict the forgotten class on the
$L_2$-ball of radius $R$ around $x$. Consequently the smoothed model's
recovery radius satisfies $\bar r_t(x;\theta) \geq R$.
\end{theorem}

\begin{proof}
Direct application of \cite{cohen2019certified}, Theorem 1, to the
indicator $\mathbf{1}\{m_t(\cdot;\theta) > 0\}$ on input $x$, with the
``runner-up'' class being any non-forget label.
\end{proof}

\begin{corollary}[Generalization of certified radius]
\label{cor:gen}
Let $\widehat{\mathcal L}_\sigma$ be the empirical Monte-Carlo estimator of
\eqref{eq:smoothed} on $n_f$ forget samples with $m$ noise draws each. Under
standard Rademacher complexity assumptions on the margin-function class
$\mathcal F$,
\[
\bigl|\mathcal L_\sigma(\theta) - \widehat{\mathcal L}_\sigma(\theta)\bigr|
\;\leq\; 2 \mathfrak R_{n_f m}(\mathcal F) + B \sqrt{\tfrac{\log(2/\delta)}{2 n_f m}}
\quad\text{w.p.\ } \geq 1 - \delta.
\]
Combined with Theorem~\ref{thm:smooth}, the certified radius $R$ generalizes:
the population fraction of forget inputs with smoothed recovery radius at
least $R$ is at least the empirical fraction minus the right-hand side above.
\end{corollary}

\paragraph{Why this is the missing piece.} The main paper's eq.~(2) is an
\emph{empirical} adversarial unlearning objective in the AMUN
\cite{krishna2025amun} mold: minimize a worst-case inner attack. It carries
no certificate. Equation~\eqref{eq:smoothed} above is the first
machine-unlearning objective that yields a per-input certified lower bound
on recovery radius, and Corollary~\ref{cor:gen} transports the per-input
guarantee to the forget distribution at standard PAC rates. The existing
audit machinery in this repository provides the matched empirical
counterpart: a predicted-vs-realized radius plot is now a direct
falsification test of \eqref{eq:smoothed}'s training.

\subsection{Empirical claims, restated as theorem checks}
\label{sec:restated}

The four theorems above re-cast the benchmark results in §4--§6 of the main
paper:

\begin{enumerate}[leftmargin=1.5em,topsep=0pt]
  \item \textbf{Recovery--certification scatter (Thm 1).} The
    cross-method ($\hat L, \mathrm{median}(\hat r)$) scatter plot directly
    tests \eqref{eq:thm1-pointwise}. Methods on or above the line
    $\hat r = -\bar m / \hat L$ are tight; methods strictly above it have
    margin headroom (an \emph{actionable} signal that more aggressive
    unlearning is feasible without crossing the certificate).
  \item \textbf{Selective vs.\ generic (Thm 2).} The forget--retain $\hat
    p_F, \hat p_R$ table per method becomes a Pinsker-bounded KL test; the
    SCRUB row formally certifies as \emph{vacuous} selective leakage at the
    audited scale; the ORBIT/CE-U/BalDRO rows pass with a quantitative
    variational gap.
  \item \textbf{Local-delivery negative (Thm 3).} The mixed-surface
    stage-2 failure is not a bug-report — it is the only possible outcome
    under the covering assumption (Lemma~\ref{lem:cover}). The paper should
    stop offering ``mixed-surface stage-2 as future work''; it is
    \emph{provably} the wrong target.
  \item \textbf{Certified objective (Thm 4).} The proposed objective
    \eqref{eq:smoothed} replaces the inner attack in eq.~(2) of the main
    paper. A single $\sigma$-sweep on ORBIT or CE-U produces a
    \emph{predicted}-radius curve from Theorem~\ref{thm:smooth}; the
    existing audit produces a matched \emph{measured} curve. Their gap is
    the new headline figure.
\end{enumerate}

\subsection{What this changes in the manuscript}

These results restructure the paper's relationship to the literature:

\begin{itemize}[leftmargin=1.5em,topsep=0pt]
  \item RURK \cite{hsu2025rurk} establishes high-dimensional
    \emph{inevitability} of residual knowledge in expectation; Theorem~1
    converts inevitability into a \emph{quantitative budget} the audit can
    estimate without an oracle.
  \item \cite{yu2025impossibility} establishes path-dependence
    impossibility for retrain equivalence in multi-stage training;
    Theorem~3 establishes a \emph{different} impossibility — multi-surface
    worst-case invariance — for the same end-goal (distinguishability),
    but with a topological rather than path-dependent mechanism.
  \item AMUN \cite{krishna2025amun} uses adversarial examples as an
    \emph{unlearning} primitive; Theorem~4 inverts the role and uses
    smoothing as a \emph{certified-defense} primitive against post-hoc
    recovery attacks.
\end{itemize}

The benchmark is no longer the headline; it becomes the experimental
validation of a small theory.

% ---------------------------------------------------------------------
% Bibliography hooks added if not already in main paper:
%
% \bibitem{cohen2019certified}
%   Jeremy Cohen, Elan Rosenfeld, J.~Zico Kolter.
%   Certified adversarial robustness via randomized smoothing.
%   ICML 2019.
%
% \bibitem{yu2025impossibility}
%   Yu et al. On the impossibility of retrain equivalence in machine
%   unlearning. arXiv:2510.16629, 2025.
%
% \bibitem{krishna2025amun}
%   Krishna et al. AMUN: Adversarial machine UNlearning.
%   arXiv:2503.00917, 2025.
% ---------------------------------------------------------------------
